% Discussion, limitations, threats to validity, conclusion, artifact availability.
\section{Discussion}
\label{sec:discussion}

\paragraph{What band share implies for warm-start design} A warm start is usually judged by how close it is to the optimum: its overlap, or the quality of the classical solution it encodes. Our results say that this is the wrong objective. At fixed overlap, initial mass on near-optimal competitors \emph{lowers} ground-state sampling (coefficient \BSCoefBand\ on $\log_{10}m_5$). QAOA amplifies the low-energy band as a whole, so the ground state receives the share of that amplification it already holds.

A good warm start therefore concentrates the low-energy band on the optimum, rather than simply placing weight near it. A relaxation that confidently commits a few marginal decisions the wrong way does the opposite: it builds a band of cheap near-misses. This reasoning uses nothing specific to unit commitment, only a diagonal cost Hamiltonian and a product-state warm start. We expect it to apply to other QUBO-encoded problems, but we have not tested that.

\paragraph{A practical decision rule} Within its training distribution, the learned start is the better choice. It samples the ground state more reliably and is insensitive to the regularisation. Outside it, the relaxation is. Below the training demand range the learned start falls under the LP, and on new profile shapes it matches it. The LP has no training distribution to fall out of. That is a genuine advantage, and it argues for falling back to the relaxation whenever an instance looks unlike the training data.

\paragraph{Why this matters without a quantum advantage} Simulated annealing solves these instances in milliseconds, so nothing here is an argument for running unit commitment on a quantum computer today. The contribution is a statement about QAOA itself: which property of the input state decides what the algorithm returns. It is a property of the state and of the algorithm's amplification, not of any device. It therefore informs warm-start design for whatever hardware eventually runs QAOA at scale, and it is falsifiable in simulation now.

\paragraph{Limitations}
\begin{itemize}
\item \textbf{One system.} We studied one stylised fleet, with cost coefficients that are assumptions rather than measured data.
\item \textbf{Simplified commitment.} Commitment is fixed-block: a committed unit runs at full output, with no economic dispatch. Minimum up and down times are not enforced at all; the encoding offers only a quadratic relaxation, which was switched off, and the exact constraint would need ancillas.
\item \textbf{Exact simulation.} Circuits were simulated exactly and noise-free. Only finite-shot expectation estimates were modelled; gate and decoherence noise were not.
\item \textbf{Depth and size.} Depth was $p\leq3$, and only $p=1$ at \NLfour\ qubits, the largest size studied. The \NLfour-qubit instances cover \LfourCoverage\% of the demand distribution; we report reweighted estimates, but they apply only to that share.
\item \textbf{Band share is diagnostic, not predictive.} It needs the full spectrum, so it explains why a warm start works but cannot select one at scale.
\item \textbf{The learned start's errors grow.} Its wrong pins grow with problem size, \SlopeRatioMin--\SlopeRatioMax\ times more slowly than the LP's, but they do grow.
\end{itemize}

\paragraph{Threats to validity} These go beyond the limitations above.
\begin{itemize}
\item \textbf{Optimiser.} One stopping rule and one parameter initialisation were used for every arm. A different optimiser could reorder the arms; Nelder--Mead did not.
\item \textbf{Choosing $\varepsilon$.} Choosing $\varepsilon$ on the evaluation data favours noisy arms, so cross-validated numbers are primary.
\item \textbf{Clustering.} Scenarios cluster by structure, so every headline comparison is repeated with one value per structural group, where the gap between the LP and C-milp at the LP's in-sample best $\varepsilon$ is borderline.
\item \textbf{Fitted exponent.} The band-share exponent was fitted on the same data it explains. Held-out arms, held-out groups and the unchanged transfer to \NLfour\ qubits guard against overfitting, but do not replace a pre-registered test.
\item \textbf{Target.} The ground state is that of the QUBO, not the operational optimum; the encoding gap (Section~\ref{sec:gap}) separates the two.
\item \textbf{Multiplicity.} Many comparisons are reported, with intervals that are not corrected for multiplicity.
\end{itemize}

\paragraph{Future work} The most useful next step is a \emph{computable proxy} for band share, such as the confidence of the learned predictor or the LP's count of near-fractional pins, which would make the diagnostic a design criterion. Beyond that: runs on hardware, where shot and gate noise act together; larger systems and deeper circuits; and other problem classes, to test whether band share generalises beyond unit commitment.

\section{Conclusion}
\label{sec:conclusion}

The quality of what warm-started QAOA returns is decided by band share, not by overlap: the ground state's share of the initial low-energy band. Band share explains ground-state sampling across seven warm-start sources, depths and two system sizes, predicts held-out arms, and transfers between sizes unchanged.

On unit commitment it explains why the LP relaxation, despite a higher approximation ratio, samples the optimum less reliably than a learned start: it pins marginal units wrongly and builds a band of near-misses. The learned start is better in distribution and fails outside it, where the relaxation, which has no distribution, remains the safe choice. A classical sampler solves all of these instances in milliseconds. The finding is about QAOA, not about an advantage.

\section*{Artifact availability}
The code, run stores, trained models and the scripts that regenerate every number, table and figure in this paper from the stored results are available in the project repository (URL to be added for the camera-ready version), together with a reproducibility guide listing seeds, package versions and regeneration steps.
